% Parallel R course: 2026 update module
% Insert this file after the existing single-node parallel section and before Rcpp.
% Requires: listings, xcolor, hyperref. Optional: booktabs.

\section{Modern parallel R workflows}

\begin{frame}{Learning objectives for the updated module}
\begin{itemize}
  \item Choose an appropriate execution model: sequential, vectorised, process-based, scheduler-based, or GPU.
  \item Use the \texttt{future} ecosystem for portable parallel code.
  \item Run reproducible parallel simulations.
  \item Scale the same computation from a workstation to a Slurm cluster.
  \item Recognise when faster I/O or an analytical database is better than adding workers.
\end{itemize}
\end{frame}

\begin{frame}{A modern performance workflow}
\begin{enumerate}
  \item Measure the real bottleneck: CPU, memory, I/O, or communication.
  \item Improve the algorithm and data representation.
  \item Use vectorised or compiled routines.
  \item Parallelise only sufficiently large, independent tasks.
  \item Scale across nodes only when one node is insufficient.
  \item Verify correctness, reproducibility, memory use, and speedup.
\end{enumerate}
\vspace{0.2cm}
\alert{Parallel execution is not automatically faster.} Startup, data transfer, serialisation, and load imbalance can dominate short tasks.
\end{frame}

\begin{frame}[fragile]{How many workers should I use?}
\begin{block}{Use the resources available to the current process}
\begin{lstlisting}[language=R]
library(parallelly)

availableCores()       # respects many scheduler/container limits
availableWorkers()     # workers visible to the current R session

detectCores()          # hardware/logical CPUs, not always usable CPUs
\end{lstlisting}
\end{block}
\begin{itemize}
  \item Logical processors are not necessarily equally fast on hybrid CPUs.
  \item On an HPC system, use the CPUs granted by the scheduler, not the whole node.
  \item Leave capacity for the operating system on an interactive workstation.
\end{itemize}
\end{frame}

\begin{frame}{Processes, threads, and oversubscription}
\begin{itemize}
  \item \textbf{R workers} are usually separate processes with separate memory.
  \item BLAS, OpenMP, \texttt{data.table}, and compiled libraries may create additional threads.
  \item Eight R workers with eight BLAS threads each can create up to 64 runnable threads.
  \item Oversubscription often reduces performance and makes timings unstable.
\end{itemize}
\vspace{0.2cm}
Practical rule: when using many R workers, initially set threaded numerical libraries to one thread per worker, benchmark, and then tune.
\end{frame}

\begin{frame}[fragile]{Benchmark elapsed time and memory}
\begin{block}{Prefer representative workloads}
\begin{lstlisting}[language=R]
library(bench)

x <- runif(1e6)
bench::mark(
  loop = {
    ans <- 0
    for (i in seq_along(x)) ans <- ans + x[[i]]
    ans
  },
  vectorised = sum(x),
  check = TRUE
)
\end{lstlisting}
\end{block}
\begin{itemize}
  \item Compare results, not just timings.
  \item Report elapsed time, memory allocation, workload size, worker count, and hardware.
\end{itemize}
\end{frame}

\subsection{The future ecosystem}

\begin{frame}{Why teach the \texttt{future} ecosystem?}
\begin{itemize}
  \item It separates \textbf{what} should be evaluated from \textbf{where and how} it runs.
  \item The computational code can remain stable while the execution plan changes.
  \item Common front ends include \texttt{future.apply}, \texttt{furrr}, and \texttt{doFuture}.
  \item Backends cover local processes, remote workers, and scheduler-managed HPC jobs.
\end{itemize}
\vspace{0.2cm}
Keep \texttt{parallel} and \texttt{foreach} in the course as foundations and legacy interfaces, but use \texttt{future} as the main portable workflow.
\end{frame}

\begin{frame}[fragile]{From sequential to parallel with one change}
\begin{block}{Portable parallel map}
\begin{lstlisting}[language=R]
library(future)
library(future.apply)

f <- function(i) {
  x <- rnorm(2e6)
  c(task = i, statistic = mean(exp(x)))
}

plan(sequential)
a <- future_lapply(1:8, f, future.seed = TRUE)

plan(multisession, workers = 4)
b <- future_lapply(1:8, f, future.seed = TRUE)

identical(a, b)
plan(sequential)
\end{lstlisting}
\end{block}
\end{frame}

\begin{frame}[fragile]{Parallel tidy workflows with \texttt{furrr}}
\begin{block}{The \texttt{purrr} interface with futures}
\begin{lstlisting}[language=R]
library(future)
library(furrr)

plan(multisession, workers = 4)

results <- future_map(
  1:20,
  ~ mean(rnorm(1e6, mean = .x)),
  .options = furrr_options(seed = 2026)
)

plan(sequential)
\end{lstlisting}
\end{block}
Use this interface when students already know \texttt{purrr::map()}.
\end{frame}

\begin{frame}[fragile]{Keep existing \texttt{foreach} code with \texttt{doFuture}}
\begin{block}{Modern backend for a familiar interface}
\begin{lstlisting}[language=R]
library(future)
library(doFuture)
library(foreach)

registerDoFuture()
plan(multisession, workers = 4)

ans <- foreach(i = 1:20,
               .combine = c,
               .options.future = list(seed = TRUE)) %dopar% {
  mean(rnorm(1e6, mean = i))
}

plan(sequential)
\end{lstlisting}
\end{block}
This provides a migration path rather than requiring an immediate rewrite of all \texttt{foreach} material.
\end{frame}

\begin{frame}{Choosing a local future plan}
\begin{itemize}
  \item \texttt{sequential}: debugging, teaching, and a reproducible baseline.
  \item \texttt{multisession}: separate R processes; portable across Windows, macOS, and Linux.
  \item \texttt{multicore}: forked processes; available on supported POSIX systems, not Windows.
  \item Explicit cluster backends: use when workers are remote or require custom startup.
\end{itemize}
\vspace{0.2cm}
For a mixed classroom, \texttt{multisession} is the safest default. Always return to \texttt{plan(sequential)} after an exercise.
\end{frame}

\subsection{Reproducibility and diagnostics}

\begin{frame}[fragile]{Reproducible parallel random numbers}
\begin{block}{Do not rely only on \texttt{set.seed()} before parallel mapping}
\begin{lstlisting}[language=R]
library(future.apply)

plan(multisession, workers = 4)
set.seed(2026)

z1 <- future_lapply(1:10, function(i) rnorm(3),
                    future.seed = TRUE)

set.seed(2026)
z2 <- future_lapply(1:10, function(i) rnorm(3),
                    future.seed = TRUE)

identical(z1, z2)
plan(sequential)
\end{lstlisting}
\end{block}
Reproducibility must hold when the worker count or backend changes.
\end{frame}

\begin{frame}{Common reasons for disappointing speedup}
\begin{itemize}
  \item Tasks are too short compared with worker startup and communication.
  \item Large global objects are copied or serialised for every task.
  \item Workers compete for memory bandwidth or disk I/O.
  \item Work is unevenly distributed among workers.
  \item Nested parallelism oversubscribes CPUs.
  \item A sequential fraction limits total speedup (Amdahl's law).
\end{itemize}
\vspace{0.2cm}
Measure serial and parallel versions using the same inputs and verify that they return equivalent results.
\end{frame}

\begin{frame}[fragile]{Chunking: amortise scheduling overhead}
\begin{block}{Prefer fewer substantial tasks to millions of tiny tasks}
\begin{lstlisting}[language=R]
library(future.apply)

plan(multisession, workers = 4)

ids <- split(1:10000, cut(1:10000, 40, labels = FALSE))
ans <- future_lapply(ids, function(block) {
  vapply(block, expensive_function, numeric(1))
})
ans <- unlist(ans, use.names = FALSE)
\end{lstlisting}
\end{block}
The best chunk size depends on task duration, variability, memory use, and backend latency.
\end{frame}

\subsection{From one node to Slurm}

\begin{frame}{A scheduler changes resource ownership}
\begin{itemize}
  \item The login node is for editing, compiling, and submitting jobs, not heavy computation.
  \item Slurm grants CPUs, memory, accelerators, nodes, and wall time to a job.
  \item R should discover and use only the resources allocated to that job.
  \item Multi-node execution adds network communication and distributed failure modes.
\end{itemize}
\end{frame}

\begin{frame}[fragile]{A minimal Slurm batch job}
\begin{block}{\texttt{run\_analysis.sbatch}}
\begin{lstlisting}[language=bash]
#!/bin/bash
#SBATCH --job-name=parallel-r
#SBATCH --nodes=1
#SBATCH --cpus-per-task=8
#SBATCH --mem=32G
#SBATCH --time=00:30:00
#SBATCH --output=logs/%x_%j.out

module load R
Rscript analysis.R
\end{lstlisting}
\end{block}
Submit with \texttt{sbatch run\_analysis.sbatch}; inspect the queue with \texttt{squeue -u \$USER}.
Adapt partition, account, module name, memory, and time to the target system.
\end{frame}

\begin{frame}[fragile]{Use Slurm through \texttt{future.batchtools}}
\begin{block}{Change the plan, not the analytical function}
\begin{lstlisting}[language=R]
library(future)
library(future.apply)
library(future.batchtools)

plan(
  future.batchtools::batchtools_slurm,
  resources = list(
    time = "00:20:00",
    mem = "4G",
    ncpus = 1
  )
)

ans <- future_lapply(files, analyse_file,
                     future.seed = TRUE)
\end{lstlisting}
\end{block}
Resource field names depend on the cluster template. Test with a small job before scaling up.
\end{frame}

\begin{frame}{Single node or multiple nodes?}
\begin{itemize}
  \item Start with one node: lower latency, simpler debugging, and usually faster data access.
  \item Use multiple nodes when memory, throughput, or turnaround requirements exceed one node.
  \item Independent simulations and file-level processing scale more easily than tightly coupled algorithms.
  \item For communication-heavy numerical algorithms, specialised MPI or distributed libraries may be more appropriate.
\end{itemize}
\end{frame}

\subsection{Data movement before parallelism}

\begin{frame}{Sometimes the bottleneck is data, not R code}
\begin{itemize}
  \item CSV parsing, repeated scans, and object copying can dominate the computation.
  \item Columnar formats such as Parquet reduce unnecessary I/O for analytical workloads.
  \item Apache Arrow supports columnar and dataset-oriented workflows.
  \item DuckDB can execute analytical SQL directly over local files and tabular data.
\end{itemize}
\vspace{0.2cm}
Before adding workers, reduce the amount of data read, copied, and returned.
\end{frame}

\begin{frame}[fragile]{Query Parquet data without loading everything}
\begin{block}{Arrow dataset with lazy filtering and aggregation}
\begin{lstlisting}[language=R]
library(arrow)
library(dplyr)

sales <- open_dataset("data/sales", format = "parquet")

summary <- sales |>
  filter(year == 2026) |>
  group_by(region) |>
  summarise(revenue = sum(amount, na.rm = TRUE)) |>
  collect()
\end{lstlisting}
\end{block}
Only call \texttt{collect()} when the reduced result is needed as an in-memory R object.
\end{frame}

\begin{frame}[fragile]{Use DuckDB for analytical queries}
\begin{block}{Push computation into an analytical database}
\begin{lstlisting}[language=R]
library(DBI)
library(duckdb)

con <- dbConnect(duckdb())
result <- dbGetQuery(con, "
  SELECT region, AVG(value) AS mean_value
  FROM read_parquet('data/*.parquet')
  GROUP BY region
")
dbDisconnect(con, shutdown = TRUE)
\end{lstlisting}
\end{block}
Compare this approach with manually splitting files among R workers.
\end{frame}

\subsection{Compiled and accelerator execution}

\begin{frame}{Four distinct acceleration strategies}
\begin{itemize}
  \item \textbf{Vectorisation}: use optimised compiled routines already available to R.
  \item \textbf{Process parallelism}: distribute independent R tasks among workers.
  \item \textbf{Compiled extensions}: move a hot loop to C++ with \texttt{Rcpp}.
  \item \textbf{Accelerators}: use GPU-aware libraries for suitable dense, batched, or model-training workloads.
\end{itemize}
These strategies can be combined, but nested parallelism and data transfers must be controlled.
\end{frame}

\begin{frame}{Extend the Mandelbrot example with \texttt{RcppParallel}}
\begin{itemize}
  \item Keep the existing three implementations: R, R with a C++ kernel, and complete C++.
  \item Add a fourth implementation using \texttt{RcppParallel::parallelFor}.
  \item Partition grid indices among threads and write each output element once.
  \item Compare elapsed time and scaling for 1, 2, 4, and 8 threads.
  \item Discuss grain size, memory bandwidth, thread safety, and reproducibility.
\end{itemize}
\end{frame}

\begin{frame}[fragile]{GPU tensors from R with \texttt{torch}}
\begin{block}{Use an accelerator-aware library rather than explicit GPU loops}
\begin{lstlisting}[language=R]
library(torch)

device <- if (cuda_is_available()) torch_device("cuda") else torch_device("cpu")

x <- torch_randn(c(4000, 4000), device = device)
y <- torch_randn(c(4000, 4000), device = device)
z <- torch_matmul(x, y)

z_cpu <- z$to(device = "cpu")
\end{lstlisting}
\end{block}
GPU speedup depends on operation size, supported kernels, device memory, and host--device transfer costs.
\end{frame}

\subsection{Updated practical exercise}

\begin{frame}{Capstone: parallel model evaluation}
\begin{enumerate}
  \item Create a fixed resampling or hyperparameter grid.
  \item Implement one deterministic function that evaluates a single grid row.
  \item Run it sequentially and verify results.
  \item Run locally with \texttt{future\_map()} or \texttt{future\_lapply()}.
  \item Run the same function through a Slurm backend.
  \item Report elapsed time, speedup, efficiency, memory, and failed tasks.
\end{enumerate}
\end{frame}

\begin{frame}[fragile]{Capstone skeleton}
\begin{block}{Backend-independent worker function}
\begin{lstlisting}[language=R]
library(future)
library(furrr)

grid <- expand.grid(
  penalty = 10^seq(-4, 1, length.out = 20),
  split = 1:10
)

evaluate_one <- function(i) {
  row <- grid[i, ]
  # fit on the training split; score on the test split
  data.frame(i = i, penalty = row$penalty,
             split = row$split, score = NA_real_)
}

plan(multisession, workers = 4)
scores <- future_map_dfr(
  seq_len(nrow(grid)), evaluate_one,
  .options = furrr_options(seed = 2026)
)
plan(sequential)
\end{lstlisting}
\end{block}
Replace the placeholder score with the course model and dataset.
\end{frame}

\begin{frame}{Updated summary}
\begin{itemize}
  \item Optimise the algorithm and data movement before parallelising.
  \item Use \texttt{parallel} to explain processes, PSOCK, and forked execution.
  \item Use \texttt{future} as the primary portable parallel interface.
  \item Treat reproducible random numbers and worker cleanup as core correctness topics.
  \item Use Slurm-aware backends for HPC scaling.
  \item Consider Arrow, DuckDB, compiled code, or GPUs when they match the bottleneck.
\end{itemize}
\end{frame}

\begin{frame}{References and documentation}
\scriptsize
\begin{itemize}
  \item Futureverse: \url{https://www.futureverse.org/}
  \item \texttt{future.batchtools}: \url{https://future.batchtools.futureverse.org/}
  \item \texttt{parallelly}: \url{https://parallelly.futureverse.org/}
  \item Arrow R documentation: \url{https://arrow.apache.org/docs/r/}
  \item DuckDB R client: \url{https://duckdb.org/docs/stable/clients/r}
  \item \texttt{RcppParallel}: \url{https://rcppcore.github.io/RcppParallel/}
  \item torch for R: \url{https://torch.mlverse.org/}
  \item Slurm documentation: \url{https://slurm.schedmd.com/documentation.html}
\end{itemize}

% Three additional Beamer slides: GPU + parallel R
% Insert after the Rcpp/RcppParallel material or before the course summary.
% No images or additional LaTeX packages are required.

\subsection{GPU computing with R}

\begin{frame}{CPU parallelism and GPU acceleration}
\small
\begin{itemize}
  \item A CPU has a relatively small number of powerful, general-purpose cores; a GPU provides many simpler execution units for highly parallel numerical workloads.
  \item GPU acceleration is most useful for large matrix operations, tensor computations, deep-learning models, and large batches of similar operations.
  \item CPU parallelism is usually more suitable for irregular tasks, branching algorithms, many small independent jobs, and workflows dominated by disk or network access.
  \item A GPU does not accelerate ordinary R code automatically: computations must use a GPU-enabled library and data must be transferred to GPU memory.
\end{itemize}
\vspace{0.15cm}
\alert{Practical rule:} first create and verify a CPU baseline; use a GPU only when the workload is sufficiently large to compensate for setup and data-transfer costs.
\end{frame}

\begin{frame}[fragile]{GPU tensors in R with \texttt{torch}}
\begin{block}{The same tensor operation on the available device}
\begin{lstlisting}[language=R,basicstyle=\ttfamily\scriptsize]
library(torch)

use_gpu <- cuda_is_available()
device <- if (use_gpu) torch_device("cuda") else torch_device("cpu")
cat("Using:", device$type, "\n")

n <- 4000
x <- torch_randn(n, n, device = device)
y <- torch_randn(n, n, device = device)

z <- torch_matmul(x, y)
if (use_gpu) cuda_synchronize()

result <- z$mean()$item()
rm(x, y, z)
\end{lstlisting}
\end{block}
\begin{itemize}
  \item \texttt{cuda\_is\_available()} checks whether CUDA can be used.
  \item Keep intermediate tensors on the GPU; copy back only the final result needed in R.
  \item Synchronise before timing because GPU operations may execute asynchronously.
\end{itemize}
\end{frame}

\begin{frame}[fragile]{Requesting a GPU with Slurm}
\begin{block}{Example batch script; adapt resource names to the cluster}
\begin{lstlisting}[language=bash,basicstyle=\ttfamily\scriptsize]
#!/bin/bash
#SBATCH --job-name=torch-r
#SBATCH --nodes=1
#SBATCH --ntasks=1
#SBATCH --cpus-per-task=4
#SBATCH --mem=16G
#SBATCH --time=00:15:00
#SBATCH --gres=gpu:1
#SBATCH --output=logs/%x_%j.out

module load R
# The cluster may also require CUDA or a suitable container/module.
Rscript gpu_torch.R
\end{lstlisting}
\end{block}
\begin{itemize}
  \item Request a GPU explicitly; do not run GPU exercises on a login node.
  \item The Slurm GPU option, partition name, modules, drivers, and supported CUDA version are cluster-specific.
  \item For this four-hour course, one shared demonstration GPU is enough; the main hands-on exercises can remain CPU-based.
\end{itemize}
\end{frame}

\end{frame}
